% Part 1: Introduction. 14 frames, lecture, session 0:00-0:20.
\section{Introduction}\subsection{}

% Instructor intro. Keep it to a minute.
\begin{frame}{About me}
\begin{wideitemize}
  \item \ybb{Xugan Chen} -- PhD candidate, Financial Economics, Yale SOM.
        At Yale since 2020 (predoc 2020--2022, PhD since 2022).
  \item \ybb{TA for:} Entrepreneurial Finance; Venture Capital and Private Equity;
        Statistical Foundations; Behavioral Finance.
  \item \ybb{Research:} innovation; AI in finance and economics; corporate
        finance; entrepreneurial finance.
\end{wideitemize}
\vs
\begin{center}
\yob{Feedback is welcome -- now, or any time after today.}
\end{center}
\end{frame}

% Poll 1: who has used what.
\begin{frame}{Getting to know you}
Show of hands. Which side of each pair are you here for?
\vs
\begin{center}
\begin{tabular}{rcl}
\ybb{Public markets}      & vs & \ybb{Private markets} \\[0.6em]
\ybb{Industry}            & vs & \ybb{Academia} \\[0.6em]
\ybb{Quant research}      & vs & \ybb{Fundamental research} \\[0.6em]
\ybb{Equities}            & vs & \ybb{Other asset classes} \\
\end{tabular}
\end{center}
\vs
Today's task sits in \yob{private markets}, and the workflow is the same
whichever side you raised your hand on.
\end{frame}

% Poll 2: the tool ladder. Figure adapted from Kyle Jensen.
\begin{frame}{Getting to know you: How do you use AI tools?}
\vs
\begin{center}
\tabfont
\begin{tabular}{ll}
\toprule
\textbf{Tool} & \textbf{What you actually do} \\
\midrule
Browser chat & Copy code out, paste it in, run it, paste the error back \\
Agentic IDE (Cursor)              & Tell it what to do; it edits files inside the editor \\
\rowcolor{yaleblue!10} \ybb{Terminal agent} (Claude Code, Codex) & \ybb{Full environment}: reads, writes, runs commands, plans \\
Terminal agent plus MCP tools            & Wire it to databases, APIs, other services \\
Agent in a container              & One instruction, and leave alone \\
\bottomrule
\end{tabular}
\end{center}
\vs
\ybb{Today is the terminal agent}, with one look at agents in a container at the end.
Agent in a container is not better; it is more setup.
\source{Ladder adapted from Kyle Jensen.}
\end{frame}

% What the tool was a year ago: chat, copy, paste.
\begin{frame}{Working with AI, one year ago}
\begin{wideitemize}
  \item A \yob{chat window}. We describe the problem in prose, it hands back a
        snippet, we paste the snippet.
  \item It could not see our data, our files, or our errors. We were the
        clipboard between the model and the project.
  \item The code was \yob{roughly right and often not runnable}, which meant we
        still had to be the one who ran it.
\end{wideitemize}
\vs
\begin{center}
That is the version most people still have in mind.
\end{center}
\end{frame}

% What it is now: it acts inside your project.
\begin{frame}{Working with AI, now: it acts in our project}
An agent that lives in our terminal, inside our project directory.
\vs
\begin{wideitemize}
  \item \ybb{Reads and writes files} in the repo -- our data, our code, our notes.
  \item \ybb{Runs commands} -- installs packages, executes scripts, reads the errors.
  \item \ybb{Iterates} -- sees its own output fail and tries again, without us retyping.
\end{wideitemize}
\takeaway{It does not hand us a snippet to paste. It acts in the project and \ybb{lives with the consequences.}}
\end{frame}

% The instruct-and-review loop, as a diagram.
\begin{frame}{The loop is easy. Instruct and review is the job.}
\vs
\begin{center}
\begin{tikzpicture}[node distance=6mm]
  \node[flowbox_us]                (a) {instruct};
  \node[flowbox, right=of a]        (b) {it proposes\\and acts};
  \node[flowbox_us, right=of b]    (c) {we review};
  \node[flowbox, right=of c]        (d) {iterate};
  \draw[flowarrow] (a) -- (b);
  \draw[flowarrow] (b) -- (c);
  \draw[flowarrow] (c) -- (d);
  \draw[flowarrow] (d.south) -- ++(0,-0.7) -| (a.south);
  \node[flowlabel, above=1mm of a] {Human};
  \node[flowlabel, above=1mm of c] {Human};
\end{tikzpicture}
\end{center}
\vs
\begin{wideitemize}
  \item The two boxes with our names on them are \yob{the entire job}.
  \item Everything we add today -- planning, verification, generalization --
        makes the instruction sharper, or makes review possible on output
        that is \ybb{too large to read}.
\end{wideitemize}
\takeaway{For example, today we will sign off on a dataset of 37,197 rows without reading one of them.}
\end{frame}

% Why review is the hard half: it checks its own work.
\begin{frame}{Why reviewing is hard: it checks its own work}
\begin{wideitemize}
  \item Ask it to build a dataset and you usually get a check you did not ask
        for: \ylb{``I reconciled the total to the reported figure -- it matches.''}
  \item That is genuinely useful. It is also the trap, because \ybb{a check looks
        at one thing}, and the agent decided which thing.
\end{wideitemize}
\vs
\begin{center}
\begin{tikzpicture}[node distance=8mm]
  \node[flowbox, text width=3.4cm]                    (t) {the total number\\matches exactly};
  \node[flowbox_warn, right=of t, text width=3.4cm]   (n) {BUT every name\\in it is wrong};
  \node[flowbox_warn, right=of n, text width=3.4cm,
        font=\small\bfseries]                         (z) {the check never\\looked at names};
  \draw[flowarrow] (t) -- (n);
  \draw[flowarrow, color=yo!70] (n) -- (z);
\end{tikzpicture}
\end{center}
\vs
\takeaway{A check that passes tells you about \ybb{the part it looked at}. It says nothing about the rest.}
\begin{center}
\yob{We are going to meet this four times this afternoon}, and it looks different every time.
\end{center}
\end{frame}

% The objective for the afternoon.
\begin{frame}{Today's objective}
\begin{wideitemize}
  \item \yob{Not} a feature walkthrough -- the menu of options will have changed by October.
  \item \yob{Not} an intro to programming -- most of you can already code.
  \item \ybb{Is} \ybb{yours to build}, with the agent, on your own machine -- not a demo you watch.
  \item \ybb{Is} practice at \ybb{checking output you cannot read} -- tens of thousands of rows.
  \item \ybb{Is} practice at \ybb{patterns that move} to the rest of your work, not just to this task.
\end{wideitemize}
\takeaway{It is never \yob{``can it write the code''}. It is \yob{``how do I know this is right''} and \yob{``how do I scale it''}.}
\end{frame}

% The roadmap: four parts.
\begin{frame}{The roadmap}
\vs
\begin{center}
\tabfont
\renewcommand{\arraystretch}{1.4}
\begin{tabular}{p{0.60\textwidth}p{0.26\textwidth}}
\toprule
\textbf{Block} & \textbf{Mode} \\
\midrule
\textbf{1.} Introduction & lecture \\
\textbf{2.} Applied AM project: a private credit dataset from SEC filings & demo, then hands-on \\
\rowcolor{yaleblue!10} \quad \ybb{Break}, 10 minutes -- halfway through the project & \\
\quad \textit{...the project continues}: generalize, then analyse & hands-on \\
\textbf{3.} More useful concepts and tips & lecture \\
\textbf{4.} Advanced usage examples & demo \\
\bottomrule
\end{tabular}
\end{center}
\vs
\ybb{Half the session (in Part 2) is hands-on}. So have your
laptop out and Claude Code installed. Work on your own or with the person
next to you -- either is fine.
\end{frame}

% Install and access. 
\begin{frame}{Basic setup}
\begin{itemize}
  \item \ybb{Install} it, either way:
  \begin{itemize}
      \item With npm, if you have Node.js: \file{npm install -g @anthropic-ai/claude-code}
      \item Or the standalone installer, for Mac, Linux and Windows:
            \yb{\url{code.claude.com/docs}}
  \end{itemize}
  \item \ybb{Then} \file{cd} into a project folder and type \file{claude}. That is all.
\end{itemize}
\vspace{-0.2em}
\begin{center}
\includegraphics[width=0.86\textwidth]{claude_code.png}
\end{center}
\vspace{-0.3em}
{\small \yob{If your terminal doesn't look like this, fix it now}. The rest of
the session assumes a working install.}
\end{frame}

% Tokens: the unit everything is measured in.
\begin{frame}{Tokens: the unit everything is measured in}
\begin{wideitemize}
  \item \ybb{Tokens} are the pieces text is chopped into before the model sees it.
  \begin{itemize}
      \item The short sentence ``I love cats'' is broken into \ylbb{3 tokens} (I, love, cats).
      \item A longer word like ``unbreakable'' may be split into \ylbb{2 tokens} (un and breakable).
      \item Rule of thumb: in English, \ybb{1 token is about 4 characters, or 0.75 words}.
  \end{itemize}
  \item Claude Fable/Opus/Sonnet: a \ybb{1M-token context window}, roughly
        750,000 words.
  \item Later we will work on SEC filings -- \ylbb{one filing runs 10 to 25 MB of HTML}.
  \begin{itemize}
      \item At 25 MB that is roughly \yob{6 million tokens -- six full context windows}.
      \item \ybb{So nothing ever reads the whole document.} It gets searched, sliced, and read a slice at a time. You will watch that happen this afternoon.
  \end{itemize}
\end{wideitemize}
\end{frame}

% Context window: long sessions degrade, so plan for it.
\begin{frame}{Context window: Long sessions degrade, so plan for it.}
\vspace{-0.2em}
\begin{columns}[T]
\begin{column}{0.4\textwidth}
\includegraphics[width=1\textwidth]{context_window2.png}
\end{column}
\begin{column}{0.6\textwidth}
\vs
\begin{itemize}\small
  \item A ``chat'' is \ybb{one long document} passed back and forth: system
        prompt, tool definitions, thinking, tool calls, results. All of it
        counts against the window, and we only see the output.
  \item \ybb{Write state to files.} A decision that lives only in the
        conversation is gone. Files persist.
  \item \ybb{Work in chunks.} Five to ten turns with one objective, then start
        fresh. Your agent at turn 30 is noticeably worse than at turn 3.
  \item \ybb{Compact.} \file{/compact} summarizes, writes to files, and continues.
\end{itemize}
\end{column}
\end{columns}
\source{Source: A Series on Claude Code by Paul Goldsmith-Pinkham.}
\end{frame}

% Housekeeping and the break.
\begin{frame}{Housekeeping}
\begin{wideitemize}
  \item \ybb{No homework, no grading.} Nothing is submitted.
  \item \ybb{Everything ships:} every prompt as it was run, every plan -- and the failures.
  \item \yob{No code is provided} -- We will build everything with the agent.
\end{wideitemize}
\vs
\begin{center}
\yb{\url{github.com/xuganchen/BootCamp-Introduction-to-Claude-Code}}
\end{center}
\end{frame}
